% U8 WS4 -- Henderson-Hasselbalch, buffer action and capacity. Blocks 7-8.
\documentclass{apchem-worksheet}

\apcunit{8}
\apcunittitle{Acids and Bases}
\apcdoctitle{Worksheet 4 \textbullet\ Buffers and Henderson--Hasselbalch}
\apcsubtitle{CED 8.8--8.10 \textbullet\ stoichiometry first, then Henderson--Hasselbalch}

\begin{document}
\wsheader[p$K_a$: \ce{HF} 3.14 \textbullet\ \ce{HNO2} 3.40 \textbullet\
\ce{HC2H3O2} 4.74 \textbullet\ \ce{HOCl} 7.46 \textbullet\ \ce{HCN} 9.21
\textbullet\ \ce{NH4+} 9.25. $K_a(\ce{HC2H3O2}) = 1.8\times10^{-5}$,
$K_w = 1.0\times10^{-14}$.]

\problem[]{\ced{8.9} Derive the Henderson--Hasselbalch equation from the
$K_a$ expression, showing each step.
\answerlines[4]{$K_a = [\ce{H+}][\ce{A^-}]/[\ce{HA}]$, so
$[\ce{H+}] = K_a[\ce{HA}]/[\ce{A^-}]$. Taking $-\log$ of both sides:
$-\log[\ce{H+}] = -\log K_a - \log([\ce{HA}]/[\ce{A^-}])$, that is
pH $=$ p$K_a - \log([\ce{HA}]/[\ce{A^-}])$. Inverting the fraction reverses
the sign of its logarithm, giving
pH $=$ p$K_a + \log([\ce{A^-}]/[\ce{HA}])$.}}

\problem[]{\ced{8.9} Calculate the pH of each buffer:}
\begin{parts}
  \item \SI{0.20}{\Molar} \ce{HC2H3O2} $+$ \SI{0.20}{\Molar}
        \ce{NaC2H3O2}: \answerblank[2.4cm]{4.74}
  \item \SI{0.10}{\Molar} \ce{HC2H3O2} $+$ \SI{0.30}{\Molar}
        \ce{NaC2H3O2}: \workspace[1.8cm]
        \keyonly{\begin{apcanswer}pH $= 4.74 + \log(0.30/0.10) =
        4.74 + 0.48 = \mathbf{5.22}$\end{apcanswer}}
  \item \SI{0.50}{\Molar} \ce{HF} $+$ \SI{0.25}{\Molar} \ce{NaF}:
        \answerblank[2.4cm]{2.84}
  \item \SI{0.50}{\Molar} \ce{NH4Cl} $+$ \SI{0.25}{\Molar} \ce{NH3}:
        \workspace[1.8cm]
        \keyonly{\begin{apcanswer}pH $= 9.25 + \log(0.25/0.50) =
        9.25 - 0.30 = \mathbf{8.95}$\end{apcanswer}}
\end{parts}

\problem[]{\ced{8.9} For each buffer in problem 2, state without
calculating whether its pH is above, below, or equal to p$K_a$, and why.
\answerlines[3]{(a) equal --- the ratio is 1. (b) above --- more base than
acid. (c) below --- more acid than base. (d) below --- more acid than base.
The sign of the log term is decided entirely by which member of the pair is
in excess.}}

\problem[]{\ced{8.10} \SI{1.00}{\litre} of a buffer is \SI{0.50}{\Molar} in
\ce{HC2H3O2} and \SI{0.50}{\Molar} in \ce{NaC2H3O2}.}
\begin{parts}
  \item Initial pH: \answerblank[2.2cm]{4.74}
  \item Add \SI{0.010}{\mole} \ce{NaOH}. Write the reaction and the new
        amounts. \answerlines[2]{\ce{OH^- + HC2H3O2 -> C2H3O2^- + H2O};
        \ce{HC2H3O2} $= 0.49$ mol, \ce{C2H3O2^-} $= 0.51$ mol.}
  \item New pH: \workspace[1.8cm]
        \keyonly{\begin{apcanswer}pH $= 4.74 + \log(0.51/0.49) =
        4.74 + 0.02 = \mathbf{4.76}$\end{apcanswer}}
  \item Instead add \SI{0.010}{\mole} \ce{HCl}. New pH:
        \answerblank[2.2cm]{4.72}
  \item Add \SI{0.010}{\mole} \ce{NaOH} to \SI{1.00}{\litre} of pure water
        instead. pH: \answerblank[2.2cm]{12.00}
  \item Compare the pH changes in (c) and (e) and state the point.
        \answerlines[2]{The buffer moved 0.02 units; pure water moved 5.00.
        The buffer converts added strong base into a weak conjugate base
        instead of letting $[\ce{OH-}]$ accumulate.}
\end{parts}

\problem[]{\ced{8.10} State the two-step method and say what makes each step
different.
\answerlines[4]{Step 1 --- stoichiometry: the strong reagent reacts to
completion, converting one member of the conjugate pair into the other.
Work in moles; no equilibrium constant appears. Step 2 --- equilibrium:
substitute the new amounts into Henderson--Hasselbalch. The first is a
limiting-reactant calculation, the second an equilibrium calculation, and
they must not be combined.}}

\problem[]{\ced{8.10} A student writes: ``Adding \SI{0.010}{\mole}
\ce{NaOH} to the acetate buffer, I put 0.010 into the ICE table as the
initial $[\ce{OH-}]$ and solved for $x$.'' Identify the error and describe
the correct first move.
\answerlines[3]{Strong base does not sit at equilibrium with a weak acid ---
it consumes it completely, so \ce{OH-} never appears in an ICE table. The
correct first move is a stoichiometry step subtracting \SI{0.010}{\mole}
from the weak acid and adding it to the conjugate base, after which
Henderson--Hasselbalch uses the new amounts.}}

\problem[]{\ced{8.10} Buffering capacity. Two acetate buffers each receive
\SI{0.01}{\mole} of \ce{H+}:}
\begin{parts}
  \item Buffer A is \SI{1.00}{\Molar} in each component. Ratio before and
        after: \answerblank[5cm]{1.00, then 0.98}
  \item Buffer B is \SI{1.00}{\Molar} in \ce{C2H3O2^-} and
        \SI{0.01}{\Molar} in \ce{HC2H3O2}. Ratio before and after:
        \answerblank[5cm]{100, then 49.5}
  \item Which resists better, and by roughly what factor in percentage
        change of the ratio? \answerlines[2]{A. Its ratio changed 2.0\%
        against B's 50.5\% --- about 25 times better --- even though both
        received the same quantity of acid.}
  \item State the general condition for optimal buffering and the useful
        range. \answerblank[7.9cm]{equal amounts, pH $=$ p$K_a$;
        p$K_a \pm 1$}
\end{parts}

\problem[]{\ced{8.10} Buffer design. For each target pH, choose the best
acid from the data strip and give the required
$[\text{base}]/[\text{acid}]$ ratio:}
\begin{parts}
  \item pH 3.50: \answerblank[6cm]{\ce{HNO2}, ratio 1.3}
  \item pH 7.50: \answerblank[6cm]{\ce{HOCl}, ratio 1.1}
  \item pH 9.00: \answerblank[6cm]{\ce{HCN}, ratio 0.62}
  \item Explain why \ce{HCN} would be a poor choice for a pH 3.50 buffer
        even though a ratio can be calculated.
        \answerlines[3]{p$K_a(\ce{HCN}) = 9.21$ is 5.7 units away, so the
        required ratio is about $10^{-5.7}$. Essentially no \ce{CN-} would
        be present, leaving nothing to neutralize added acid --- the
        solution would sit near pH 3.50 with no capacity at all.}
\end{parts}

\problem[]{\ced{8.8} Two acetate buffers: X is \SI{1.0}{\Molar} in each
component, Y is \SI{0.010}{\Molar} in each.}
\begin{parts}
  \item Which has the higher pH? \answerblank[5cm]{neither --- both 4.74}
  \item Which is the better buffer, and why?
        \answerlines[3]{X. The pH depends only on the ratio, which is 1 in
        both, but capacity depends on the absolute amounts. X holds 100
        times as much of each component, so it absorbs far more added acid
        or base before the ratio moves appreciably.}
\end{parts}

\begin{spiralstrip}{Unit 8 \textbullet\ blocks 1--6}
  \item pH of \SI{0.30}{\Molar} \ce{NaF}: \answerblank[2cm]{8.31}
  \item $K_b$ for acetate: \answerblank[3.2cm]{$5.6\times10^{-10}$}
  \item Lower p$K_a$ means the acid is: \answerblank[2.2cm]{stronger}
  \item Is \ce{HCl} $+$ \ce{NaCl} a buffer? \answerblank[1.6cm]{no}
  \item \ce{HF} is the weakest hydrohalic acid because of:
        \answerblank[3cm]{bond strength}
\end{spiralstrip}

\end{document}
